% F-01 总体建模流程图
\documentclass[border=5pt]{article}
\usepackage[margin=0pt,paperwidth=19cm,paperheight=11cm]{geometry}
\pagestyle{empty}
\usepackage{fontspec}
\usepackage{ctex}
\newcommand{\fontdir}{../../../00_知识库/论文写作/LaTeX模板_2026国赛版/fonts/}
\setCJKfamilyfont{song}[
  Path = \fontdir,
  UprightFont = SourceHanSerifCN-Regular.otf,
  BoldFont   = SourceHanSerifCN-Bold.otf,
  AutoFakeBold = true,
]{SourceHanSerifCN}
\newcommand*{\song}{\CJKfamily{song}}

\usepackage{tikz}
\usetikzlibrary{arrows.meta,positioning,calc,shapes.geometric,backgrounds,fit}

\definecolor{primary}{HTML}{1F3A93}
\definecolor{primaryDark}{HTML}{1A3A6E}
\definecolor{secondary}{HTML}{D9E4F2}
\definecolor{accent}{HTML}{D35400}
\definecolor{accentLight}{HTML}{FDEBD0}
\definecolor{decisionFill}{HTML}{E8F8F5}
\definecolor{decisionBorder}{HTML}{117A65}
\definecolor{gray}{HTML}{95A5A6}
\definecolor{textDark}{HTML}{333333}
\definecolor{arrowColor}{HTML}{5D6D7E}
\definecolor{groupBg}{HTML}{F4F6F7}

\tikzset{
  process/.style = {rectangle,draw=primaryDark,line width=0.8pt,fill=primary,text=white,
    minimum width=2.8cm,minimum height=0.9cm,align=center,font=\song\footnotesize,
    inner sep=4pt,rounded corners=2pt},
  processW/.style = {process,minimum width=3.6cm,minimum height=1.3cm,
    font=\song\small},
  processX/.style = {process,minimum width=17cm,minimum height=0.85cm},
  subprocess/.style = {rectangle,draw=primaryDark,line width=0.7pt,fill=secondary,
    text=textDark,minimum width=2.8cm,minimum height=0.85cm,align=center,
    font=\song\footnotesize,inner sep=4pt,rounded corners=2pt},
  subprocessW/.style = {subprocess,minimum width=17cm,minimum height=0.85cm},
  startstop/.style = {rectangle,draw=accent,line width=1.0pt,fill=accentLight,
    text=accent!85!black,minimum width=2.8cm,minimum height=0.9cm,align=center,
    font=\song\footnotesize,inner sep=4pt,rounded corners=6pt},
  arrow/.style = {-{Stealth[length=2.2mm,width=2.0mm]},line width=0.9pt,draw=arrowColor},
  arrowBold/.style = {-{Stealth[length=2.5mm,width=2.2mm]},line width=1.2pt,draw=accent},
  arrowBi/.style = {{Stealth[length=2.0mm,width=1.8mm]}-{Stealth[length=2.0mm,width=1.8mm]},
    line width=1.0pt,draw=accent},
}

\begin{document}
\begin{tikzpicture}

  % === 顶部五个阶段 — 整体右移0.5cm ===
  \node[subprocess]       (S1) at (2.0, 0) {赛题分析\\附件1/2、附录2/3/4};
  \node[process]          (S2) at (5.5, 0) {模型建立\\PDE + 边界条件};
  \node[process]          (S3) at (9.0, 0) {数值求解\\FVM + 全隐式\\+ Picard 迭代};
  \node[process]          (S4) at (12.5,0) {结果验证\\收敛性/灵敏度};
  \node[startstop]        (S5) at (16.0,0) {论文呈现\\结果 + 检验 + 评价};

  \draw[arrow] (S1) -- (S2); \draw[arrow] (S2) -- (S3);
  \draw[arrow] (S3) -- (S4); \draw[arrow] (S4) -- (S5);

  % === 中部两个核心方程 ===
  \node[processW] (EqT) at (3.5, -2.6) {
    热传导方程\\[3pt]
    $\displaystyle\rho c_p\frac{\partial T}{\partial t}
    =\frac{1}{r}\frac{\partial}{\partial r}\bigl(kr\frac{\partial T}{\partial r}\bigr)$
  };
  \node[processW] (EqC) at (11.0,-2.6) {
    Fick 扩散方程\\[3pt]
    $\displaystyle\frac{\partial C}{\partial t}
    =\frac{1}{r}\frac{\partial}{\partial r}\bigl(Dr\frac{\partial C}{\partial r}\bigr)$
  };

  \draw[arrow] (S2) -- (EqT);
  \draw[arrow] (S3) -- (EqC);
  \draw[arrowBi] (EqT) -- (EqC) node[midway,above,font=\song\footnotesize] {双向耦合};

  % === 四问递进长条 ===
  \node[startstop,minimum width=16cm,minimum height=0.85cm,
        fill=accentLight!60,draw=accent!50,text=textDark]
    (Chain) at (8.0, -4.5) {
    \textbf{同一PDE框架，逐问加码：}~~Q1定物性(附录2)~$\rightarrow$~Q2变物性(附录3)~
    $\rightarrow$~Q3长时程+终止判据~$\rightarrow$~Q4移动边界(附录4+附件2)
  };
  \draw[arrow] (EqT.south) -- ++(0,-0.5) -| (Chain.north);
  \draw[arrow] (EqC.south) -- ++(0,-0.5) -| (Chain.north);

  % === 数值方法 ===
  \node[subprocessW,minimum width=16cm] (Method) at (8.0, -6.0) {
    数值方法：~~守恒型有限体积法~~|~~全隐式时间推进~~|~~
    Picard $\omega=0.7$, tol=$10^{-10}$~~|~~Q4: 物料坐标 $\xi=r/R(t)$
  };
  \draw[arrow] (Chain) -- (Method);

\end{tikzpicture}
\end{document}